\newpage
\section{\textbf{Conclusões}}

O estudo foi desenvolvido para a SE Campus do Pici com um transformador energizado e outro em reserva, mantendo-se os níveis de curto-circuito fornecidos pelo professor e incorporando-se os disjuntores de alta tensão dos transformadores à configuração proposta. Os critérios do roteiro foram relacionados aos parâmetros da apresentação sobre proteção de subestações de 69--13,8~kV, sem transferência dos dados de potência, cabos ou curto-circuito do exemplo para a instalação estudada.

Na revisão dos TCs de alta tensão, foram adotados 1200-5 na entrada 12L1 e 800-5 em 12T1, 12T2 e nos núcleos exclusivos de alta da proteção diferencial. Foram mantidos 400-5 nos gerais de média tensão e 200-5 nos alimentadores e nos neutros. As novas relações de alta atenderam aos três critérios calculados, com limites de curto de 24.000 e 16.000~A, superiores a 15.683,2~A. A insuficiência anterior do critério de curto foi resolvida, tendo permanecido necessária a confirmação das classes, da excitação e da suportabilidade dos equipamentos especificados.

Na sobrecorrente, foram adotadas curvas IEC muito inversas C2, com alvos iniciais de aproximadamente 0,20 e 0,30~s para fase e neutro dos alimentadores. Os ajustes secundários foram recalculados para os novos TCs, com arredondamento ao passo de 0,01~A e atualização dos valores primários efetivos. Os dials e aditivos permitiram manter margens nominais superiores ou iguais a 0,40~s entre as funções temporizadas locais nos intervalos verificados, inclusive no patamar de múltiplo máximo. Foram utilizados 24~A nos alimentadores, 60~A nos gerais e no neutro do transformador, 60,8~A no primário e 120~A na entrada.

A função 50 de fase do primário foi ajustada acima da corrente passante refletida acrescida do fator de assimetria de 1,73, com atraso de 0,05~s. Nos gerais de média tensão e na entrada, os estágios de alta corrente foram mantidos desabilitados, conforme a configuração de referência dos slides. Nos alimentadores, os valores calculados a partir de 80\% da corrente máxima foram classificados como pré-ajustes, sem comprovação de alcance ao longo da rede. A liberação desses estágios permaneceu dependente das correntes de falta até o primeiro religador.

Para a proteção diferencial, foram adotados os parâmetros de sensibilidade e dupla inclinação apresentados nos slides, com compensação do grupo Dyn1 e utilização de um IED compatível com o casamento das correntes e a sensibilidade requerida. A análise das faltas externas foi realizada com correntes compensadas, sem atribuição de corrente residual ao lado em triângulo. A característica percentual e os pontos de operação foram definidos analiticamente e representados na Figura~\ref{fig:87_caracteristica_diferencial}, sem atribuição de ensaios ou simulações transitórias não realizados.

A seletividade integral com a concessionária não foi demonstrada para as faltas do setor de 69~kV. A atuação de fase de Pici em aproximadamente 0,134~s e a proteção de distância existente impuseram restrições que não foram resolvidas apenas pelos ajustes locais. Os tempos típicos de processamento e interrupção dos slides reforçaram a necessidade de conferir o tempo total de eliminação, em lugar de comparar somente os atrasos intencionais.

Dessa forma, foi estabelecida uma proposta de ajustes e de coordenação local, com identificação das condições de aplicação e das limitações. A conclusão da verificação da rede completa permaneceu dependente dos curtos mínimos, dos ajustes dos religadores, do histórico de desequilíbrio, das características efetivas dos TCs e dos tempos dos disjuntores. As expressões e tabelas foram complementadas pelos coordenogramas, pelos diagramas de tempo e pela característica diferencial, mantendo-se as condições de aplicação dos ajustes representados.
